% --------------------------- ABSTRACT --------------------------------
\begin{abstract}
\noindent
This paper turns a well-known benchmark-rebalancing anomaly into a directly testable investment process. OSEBX, the Oslo Stock Exchange Benchmark Index, is rebuilt semi-annually from a transparent Euronext rule book. I use only the variables in that rule book--trading status, electronic-order-book turnover, free float, sector, and lagged membership--to predict index additions and deletions before the official announcement, then trade the predicted changes through active and enhanced-index portfolios. A logistic GLM and XGBoost both rank future constituents well: out-of-sample AUROC exceeds 0.97 across 26 rebalancing events from 2010 to 2023. The harder problem is economic, not statistical. A naive persistence model has 95.1\% accuracy but produces zero trades. To make the forecast investable, I introduce a conditional posterior-probability threshold that deliberately sacrifices some accuracy to increase the number of predicted changes. The best active overlays generate significant Fama--French three-factor alpha, while the risk-controlled XGBoost variant remains significant when embedded in a 2\% tracking-error enhanced-index sleeve. The central conclusion is practical: in benchmark-event trading, the useful model is not the most accurate classifier, but the classifier whose errors create a diversified and capacity-aware signal.
\end{abstract}
